Angular Separator \(\psi\) of the TanArcanTau Family as Kernel Bridge Division  
Completing a Calculable Numerical E-Fold through Cosmological Clock Folding of Seven Consecutive Multiples that Propagate Radial Symmetry, Generative Winding, Equidistributed Altitude and Side Length, and the Residual Bulge of the Thinnest Equidistributal Transcendental Triangle while Remaining Phase-Locked to the Cosmological Clock by Operations that Preserve Transcendence

1. The Angular Separator as Kernel Bridge

Within the continuous TanArcanTau (Arcan(\(\tau\))) family the four fundamental angles

\[
\begin{align*}
\theta&=\arctan(\tau),\qquad
\alpha=\theta/2,\\
\phi&=\arctan(\pi),\qquad
\psi=\theta-\phi
\end{align*}
\]

are generated from the single transcendental circle constant \(\tau=2\pi\). The difference \(\psi\approx0.1503378808\) is the angular separator that stands between the full-turn projection \(\theta\) and the half-turn projection \(\phi\). In the architecture of the Cosmological Clock this separator functions as the Kernel Bridge Division: it is the irreducible angular interval that must be crossed in order to pass from the continuous transcendental source to every discrete residual structure that follows. All subsequent numerical content of the hierarchy—residual seed, intensity, bulge, altitude, side lengths—is obtained by operations performed upon \(\psi\).

2. Seven Consecutive Multiples and the Residual Seed

The Kernel Bridge is iterated seven times and reduced modulo \(2\pi\):

\[
7\psi=\frac{\pi}{3}+\delta,\qquad\delta\approx0.0051676144\ \mathrm{rad}.
\]

The integer seven is the minimal count that brings the separator into near-equivalence with the elementary six-fold lattice while leaving a measurable oriented residual \(\delta\). This residual is the generative seed of radial symmetry. Because the operations that produce \(\delta\) are field operations and integer multiplication applied to a transcendental angle, the seed itself remains transcendental. Transcendence is therefore preserved at the first discrete step.

3. Calculable Numerical E-Fold through the Cosmological Clock

The Cosmological Clock advances by the discrete tick lattice \(\varepsilon=10^{-9}\) and the scale factor \(a(N)=\mathrm{e}^{\varepsilon N}\). The same angular separator \(\psi\) that supplies the residual seed is already the microscopic generator of the Clock’s phase bias. Consequently every structure descended from \(7\psi\) remains phase-locked to the Clock’s numerical e-fold. The residual hierarchy does not float free of cosmic expansion; it is the geometric reading of the identical angular architecture that drives the late-time geometric torque. The e-fold is calculable because both the residual seed and the Clock’s phase are closed-form expressions in the single transcendental separator \(\psi\).

4. Propagation of Radial Symmetry as Generative Winding

Once the residual seed exists, the practical generator \(7\psi\) acts as a generative winding. Six applications close into the Order-6 radial fold—the practical hexagon and its interlocking residual triangles. Further applications densify the orbit (Orders 10, 12, 15, 30, \ldots). Because \(7\psi/2\pi\) is irrational the orbit is equidistributed; the residual that prevents exact equilateral closure is the same residual that guarantees equidistribution at finer scales. Radial symmetry therefore propagates as a generative winding whose every finite truncation still carries the measurable six-fold memory of the original Kernel Bridge.

5. Equidistribution of Altitude and Side Length

The even-index cell of the first closure is the thinnest residual near-equilateral triangle \(T\):

\[
\begin{align*}
s_{\mathrm{thin}}&=1.721623257385,\\
s_{\mathrm{long}}&=1.737195272475,\\
h_{\mathrm{res}}&=1.490969476641,\\
\operatorname{Area}(T)&=1.2988990741.
\end{align*}
\]

These metric quantities are transcendental invariants of the Arcan(\(\tau\)) family: they arise from trigonometric evaluations of transcendental multiples of \(\psi\) by operations (subtraction, squaring, square root, polynomial combination) that do not cancel transcendence. Because the orbit is equidistributed, congruent copies of the same altitude and side lengths reappear at every radial scale. The triangle is therefore equidistributal: its metric data are uniformly sampled across the hierarchy while remaining transcendental.

6. Residual Bulge of the Thinnest Equidistributal Transcendental Triangle

From the complementary residual area of the unit disk one extracts the dimensionless intensity

\[
\rho=\frac{A_{\mathrm{slice}}}{\operatorname{Area}(T)}\approx0.472886
\]

and returns it to the edges as the normal displacement

\[
\gamma_\rho(t)=\gamma_0(t)+c\cdot\rho\cdot4t(1-t)\,n.
\]

The resulting bulge \(\approx0.08512\) is the continuous geometric expression of the residual intensity. It is the winding method’s contribution to the dual structure: the same Kernel Bridge that folds into the residual hexagon also bulges the sides of the residual triangle. Both the bulge and the hexagon remain phase-locked to the Cosmological Clock because both descend from the identical separator \(\psi\).

7. Preservation of Transcendence under Phase-Lock

Every numerical quantity that appears in the cascade—\(\delta\), side lengths, altitude, area, \(\rho\), bulge—is obtained from \(\psi\) by a finite sequence of algebraic and elementary transcendental operations. None of these operations maps the original transcendental data onto an algebraic number. Consequently the entire residual hierarchy, while phase-locked to the discrete-tick expansion of the Cosmological Clock, preserves the transcendence of its Arcan(\(\tau\)) source.

8. Unified Statement

The angular separator \(\psi\) of the TanArcanTau family is the Kernel Bridge Division. Seven consecutive multiples of this bridge produce the residual seed \(\delta\) that organises radial symmetry as a generative winding. The winding equidistributes the transcendental altitude and side lengths of the thinnest residual triangle and, through the residual intensity \(\rho\), bulges the edges of that triangle. The whole construction completes a calculable numerical e-fold that remains phase-locked to the Cosmological Clock by operations that preserve transcendence. Residual bulge and residual hexagon are therefore the dual, Clock-locked, transcendence-preserving geometric expressions of a single angular separator.
